% 디지털 포렌식·침해사고대응(DFIR) 직무용 이력서. 서식과 공통 섹션은 common/ 에 있습니다.
%-------------------------
% Resume in Latex
% Original Author : Sourabh Bajaj
% Adaptation : Hyunggi Chang
% Draft content : Changyeob Lee
% License : MIT
%
% 직군별 이력서(resume_{dfir,malware,cert}_{ko,en}.tex)가 함께 쓰는 서식입니다.
%------------------------

\documentclass[letterpaper,11pt]{article}

\usepackage{latexsym}
\usepackage[empty]{fullpage}
\usepackage{titlesec}
\usepackage{marvosym}
\usepackage[usenames,dvipsnames]{color}
\usepackage{verbatim}
\usepackage{enumitem}
\usepackage[hidelinks]{hyperref}
\usepackage{fancyhdr}
\usepackage[english]{babel}
\usepackage{tabularx}
\usepackage{amsmath}
\usepackage{kotex}
\usepackage{needspace}

\pagestyle{fancy}
\fancyhf{}
\fancyfoot{}
\renewcommand{\headrulewidth}{0pt}
\renewcommand{\footrulewidth}{0pt}

\addtolength{\oddsidemargin}{-0.5in}
\addtolength{\evensidemargin}{-0.5in}
\addtolength{\textwidth}{1in}
\addtolength{\topmargin}{-0.5in}
\addtolength{\textheight}{1.0in}

\urlstyle{same}
\raggedbottom
\raggedright
\setlength{\tabcolsep}{0in}

\titleformat{\section}{
  \vspace{-4pt}\scshape\raggedright\large
}{}{0em}{}[\color{black}\titlerule \vspace{-2pt}]

\newcommand{\resumeItem}[1]{
  \item\small{{#1 \vspace{-2pt}}}
}

% 쪽 끝에서 소제목만 남지 않도록, 남은 공간이 부족하면 다음 쪽에서 시작합니다.
\newcommand{\resumeSubheading}[4]{
  \needspace{3\baselineskip}\vspace{-1pt}\item
    \begin{tabular*}{0.97\textwidth}[t]{l@{\extracolsep{\fill}}r}
      \textbf{#1} & #2 \\
      \textit{\small#3} & \textit{\small #4} \\
    \end{tabular*}\vspace{-5pt}
}

% 설명이 두 줄 이상이어도 아래 항목과 겹치지 않도록 minipage 대신 일반 문단으로 둡니다.
\newcommand{\resumeProject}[2]{
  \vspace{-1pt}\item
    \textbf{#1} \\
    {\small #2\par}\vspace{-3pt}
}

\newcommand{\resumeTalk}[2]{
  \vspace{-1pt}\item
    \textbf{#1} \\
    {\small\itshape #2\par}\vspace{-3pt}
}

\newcommand{\resumeSkill}[2]{
  \item\small{\textbf{#1}: #2}
}

\newcommand{\resumeAward}[3]{
  \item\small{\textbf{#1} / #2 / #3 \vspace{-2pt}}
}

\renewcommand{\labelitemii}{$\circ$}
\newcommand{\resumeSubHeadingListStart}{\begin{itemize}[leftmargin=*]}
\newcommand{\resumeSubHeadingListEnd}{\end{itemize}}
\newcommand{\resumeItemListStart}{\begin{itemize}}
\newcommand{\resumeItemListEnd}{\end{itemize}\vspace{-5pt}}


\begin{document}

%----------HEADING-----------------
\begin{tabular*}{\textwidth}{l@{\extracolsep{\fill}}r}
  \textbf{\href{https://info.yeopeva.me/}{\Large 이창엽 (Changyeob Lee)}} & Github: \href{https://github.com/YeoPEVA}{YeoPEVA} $|$ LinkedIn: \href{https://www.linkedin.com/in/changyeob-lee-yeopeva}{changyeob-lee-yeopeva} \\
  \href{https://info.yeopeva.me/}{https://info.yeopeva.me/} & Email: \href{mailto:lcy606@naver.com}{lcy606@naver.com} \\
  \href{https://noperfectsecurity.tistory.com/}{noperfectsecurity.tistory.com} & KST / GMT+9
\end{tabular*}


%-----------SUMMARY-----------------
\section{요약}
  \resumeSubHeadingListStart
    \resumeItem{디지털 포렌식과 침해사고 분석을 중심으로 역량을 쌓아온 정보보호 전공 학부생입니다.}
    \resumeItem{포렌식 대회 및 CTF, 다양한 경험을 바탕으로 흔적을 따라 원인을 좁혀가는 훈련을 해왔습니다.}
    \resumeItem{분석 결과를 재현 가능한 절차와 읽히는 보고서로 정리하는 것을 중요하게 생각합니다.}
  \resumeSubHeadingListEnd

%-----------SKILLS-----------------
\section{기술}
  \resumeSubHeadingListStart
    \resumeSkill{DFIR}{Windows·Linux 포렌식, 네트워크 포렌식(Wireshark, PCAP), 타임라인 재구성, 침해사고 분석}
    \resumeSkill{악성코드 분석}{정적·동적 분석, PE 구조 분석 / IDA Pro, OllyDbg, x64dbg, Detect It Easy}
    \resumeSkill{탐지 룰 · 위협 인텔리전스}{YARA / VirusTotal, ANY.RUN, Shodan, MITRE ATT\&CK}
    \resumeSkill{보안 운영}{보안 관제, 장비 관리, 정책·차단 관리}
    \resumeSkill{엔지니어링}{Python(분석·수집 자동화), Backend, DevSecOps, IaC}
    \resumeSkill{Languages}{Korean (Native), English (Intermediate), Japanese (Intermediate)}
  \resumeSubHeadingListEnd

%-----------PROJECTS-----------------
\section{대표 프로젝트 및 활동}
  \resumeSubHeadingListStart
    \resumeProject
      {악성코드 분석 \& YARA 룰 작성 · SCH 사이버보안연구센터}
      {2025.09 -- 현재 / 랜섬웨어·인포스틸러 샘플을 분석하고 탐지 룰과 자동화를 연구하고 있습니다.}
      \resumeItemListStart
        \resumeItem{PE 구조·임포트 테이블·문자열·난독화 여부를 먼저 보고 실행 흐름을 따라가며, 분석한 특징은 IoC와 YARA 룰로 정리합니다.}
      \resumeItemListEnd

    \resumeProject
      {Stellar-Agent · IaC-Rule Set \& MCP 기반 정책 감사 도구}
      {2025.03.03 -- 2025.12.18 / 팀 프로젝트. IaC Rule Set과 MCP 기반 AI로 보안 정책 감사, SBOM 및 보고서 생성을 자동화하는 도구를 만들었습니다. 2025 SCHU AI \& SW Festival SW 프로젝트 경진대회 대상을 받았고 학술대회 논문을 공저했습니다.}
      \resumeItemListStart
        \resumeItem{IaC 파트를 맡아 IaC 파일을 SBOM으로 바꾸는 도구의 CloudFormation 파서(문법 검증, 전용 태그 처리, 변수 추출)를 구현하고 스캐너 연동 일부에 참여했습니다.}
      \resumeItemListEnd

    \resumeProject
      {디지털 포렌식 · 침해대응 경진대회 분석}
      {2022.12 -- 2025.12 / 디스크 이미지와 로그를 교차 검증해 공격자의 실행·지속성 흔적을 추적하고 타임라인으로 재구성했습니다.}
      \resumeItemListStart
        \resumeItem{HCCC 2025에서 Web Server 이미지 분석을 담당해 접근 로그의 명령형 요청 8건을 웹셸 코드와 대조하고, 백도어 의심 파일을 부팅 시와 매시 실행하도록 등록한 root cron(지속성 설정)을 확인했습니다.}
        \resumeItem{디지털 범인을 찾아라 주니어 부문에서는 Windows 이미지의 레지스트리·프리패치·LNK·\$LogFile·\$UsnJrnl을 대조해 파일 은닉·삭제·이름 변경 흔적을 시간선으로 정리했습니다.}
      \resumeItemListEnd

    \resumeProject
      {금융보안아카데미 2025 사이버 위협 대응·분석 과정}
      {2025.07.08 -- 2025.11.07 / 금융권 위협 분석·대응 절차를 학습하고 DFIR 중심 CTF를 수행했으며, 과정 전체 최우수상(금융보안원장, 팀)을 받았습니다.}
      \resumeItemListStart
        \resumeItem{팀으로 국제 CTF 3건 포함 7개 대회에 참가했고, 화이트햇 콘테스트 예선에서는 포렌식 파트를 맡아 PCAP 숨김 메시지 복원 등 2문제를 직접 풀었고, 팀은 대학생부 4위로 본선에 진출했습니다.}
      \resumeItemListEnd

    \resumeProject
      {i-Keeper / MFT Parser 분석 및 개선}
      {2023.09 -- 2024.02 / Python. analyzeMFT 기반 CLI 도구에 PyQt GUI를 구현하고 JSON 출력 예외 처리를 보강했습니다.}

    \resumeProject
      {BoB - 중고거래사기방지플랫폼 구현}
      {2020.09 -- 2020.12 / 팀 프로젝트 · Python, PHP. 3개 플랫폼 게시물 수집과 MySQL·ELK 구축을 담당했고, 결과를 KCI 논문으로 확장했습니다.}
  \resumeSubHeadingListEnd

%-----------EXPERIENCE-----------------
\section{대표 경험}
  \resumeSubHeadingListStart
    \resumeSubheading
      {HSPACE FORENSIC LAB}{대한민국}
      {Member / HDFLAB}{2025.10 -- 현재}
      \resumeItemListStart
        \resumeItem{상·하반기 제공 DFIR 문제 풀이 및 분석 보고서 작성}
      \resumeItemListEnd

    \resumeSubheading
      {SCH 사이버보안연구센터}{순천향대학교}
      {Member / 악성코드, 익스플로잇 킷 분석 및 연구}{2025.09 -- 현재}
      \resumeItemListStart
        \resumeItem{악성코드와 익스플로잇 킷 분석 연구 수행}
      \resumeItemListEnd

    \resumeSubheading
      {대한민국 공군}{대한민국}
      {정보보호병 825기}{2021.04 -- 2023.01}
      \resumeItemListStart
        \resumeItem{정보보호 장비 관리, 관제, 정책·차단 관리 업무 수행}
        \resumeItem{훈련 대응을 통해 보안 운영 절차 경험}
      \resumeItemListEnd

    \resumeSubheading
      {Anti-root}{영남권}
      {팀장 및 설립자}{2017.05 -- 2023.02}
      \resumeItemListStart
        \resumeItem{영남권 정보보호 연구팀 설립·운영, 자체 세미나 및 프로젝트 주관}
        \resumeItem{허니팟, USB 복구, CloudGoat, YARA, 이메일 포렌식 등 발표}
      \resumeItemListEnd

    \resumeSubheading
      {부산소프트웨어마이스터고}{부산, 대한민국}
      {전공동아리 멘토 / 산학겸임교사 · 정보보호관리 교과목}{2023.04 -- 현재}
      \resumeItemListStart
        \resumeItem{2023 -- 2026학년도 전공동아리 멘토로 연속 임명되어 학생 프로젝트를 지도하고 있습니다.}
        \resumeItem{2026.03 -- 2026.07 산학겸임교사로 정보보호관리 교과목을 담당했습니다.}
      \resumeItemListEnd

    \resumeSubheading
      {i-Keeper}{대구가톨릭대학교}
      {CERT팀 / DEV팀 Member}{2023.09 -- 2024.07}
      \resumeItemListStart
        \resumeItem{CERT팀에서 MFT Parser를 분석·개선했고, DEV팀에서는 장비부장을 맡았습니다.}
      \resumeItemListEnd
  \resumeSubHeadingListEnd

%-----------AWARDS-----------------
\section{주요 수상}
  \resumeSubHeadingListStart
    \resumeAward{한국인터넷진흥원장상}{제3회 정보보호 교육·훈련 사이버공격·방어 시나리오 경진대회 (ATHENA 2026)}{2026.09.18}
    \resumeAward{장려상(한국포렌식학회장)}{제8--11회 디지털 범인을 찾아라 경진대회 (4회 연속)}{2022.12.15 / 2023.12.18 / 2024.12.17 / 2025.12.16}
    \resumeAward{최우수상(금융보안원장)}{금융보안아카데미 2025}{2025.11.07}
    \resumeAward{우수상(한국주택금융공사 사장)}{제1회 영남권 사이버 공격 방어 대회 대학팀 부문}{2025.10.15}
    \resumeAward{우수상(한국인터넷진흥원장)}{2025 제6회 호남 사이버보안 콘퍼런스 침해대응 경진대회}{2025.09.18}
    \resumeAward{장려상(한국정보보호학회 호남지부장)}{2024 제5회 호남사이버보안컨퍼런스 대학간 침해대응/분석 경진대회}{2024.10.02}
    \resumeAward{최우수상(공군참모총장)}{제44회 공군 정보통신 경연대회 해킹방어부문}{2022.10.26}
    \resumeAward{장려상(한국디지털포렌식학회장)}{2020년도 한국디지털포렌식학회 디지털포렌식 챌린지}{2020.11.20}
  \resumeSubHeadingListEnd


%-----------EDUCATION-----------------
\section{교육}
  \resumeSubHeadingListStart
    \resumeSubheading
      {순천향대학교}{아산, 대한민국}
      {SW융합대학 정보보호학과 전공 / 2027.02 졸업예정 / 학점 3.73/4.5 (4-1 종료 기준)}{2025.02 -- 2027.02}

    \resumeSubheading
      {금융보안원}{대한민국}
      {금융보안아카데미 3기 / 사이버 위협 대응·분석 과정}{2025.07.08 -- 2025.11.07}

    \resumeSubheading
      {대구가톨릭대학교}{경산, 대한민국}
      {소프트웨어융합대학 컴퓨터소프트웨어학부 컴퓨터공학전공 및 사이버 보안 복수 전공}{2020.03.01 -- 2025.02.11}
      \resumeItemListStart
        \resumeItem{중퇴 (편입) / 전체 평점 3.97/4.50, 전공 평점 4.04/4.50}
      \resumeItemListEnd

    \resumeSubheading
      {한국정보기술연구원}{대한민국}
      {차세대 보안리더 양성 프로그램 9기 / 디지털포렌식 트랙}{2020.07 -- 2021.03}
  \resumeSubHeadingListEnd


%-----------TALKS AND PAPERS-----------------
\section{발표 및 논문}
  \resumeSubHeadingListStart
    \resumeTalk{온라인 직거래 환경에서의 사기 피해에 대한 사전 예방 필요성 제언: 사기 판별 요소를 중심으로}{KCI 등재 논문, 경찰학연구 21(1), pp.249--272, 2021.03, DOI: 10.22816/polsci.2021.21.1.009}
      \resumeItemListStart
        \resumeItem{현장호, 임다연, 이창엽, 이진우, 이현호, 최원영}
      \resumeItemListEnd
    \resumeTalk{Stellar-Agent: IaC 룰셋 및 모델 컨텍스트 프로토콜(MCP)을 통합한 AI 에이전트 기반 하이브리드 \mbox{DevSecOps} 정책 감사 프레임워크}{2025 한국데이터사이언스학회 동계 종합학술대회 논문집, 2025.12.18, 논문번호 0032, p.118}
      \resumeItemListStart
        \resumeItem{정준영, 이창엽, 이정민, 이유식}
      \resumeItemListEnd
    \resumeTalk{중고 거래 사기 방지 플랫폼 제안}{한국정보보호학회 동계학술대회, 2020}
      \resumeItemListStart
        \resumeItem{현장호, 이현호, 임다연, 이진우, 이창엽, 오동빈, 최원영}
      \resumeItemListEnd
    \resumeTalk{국내 APT 및 기업 보안 사고에 대한 고찰}{한국정보보호학회 동계학술대회, 2019}
      \resumeItemListStart
        \resumeItem{이창엽, 이진우, 이정훈, 노무승, 박현상}
      \resumeItemListEnd
    \resumeTalk{카드해킹의 문제점과 대비방안}{한국정보보호학회 동계학술대회, 2018}
      \resumeItemListStart
        \resumeItem{이진우, 이정훈, 이창엽, 노무승, 오현수}
      \resumeItemListEnd
    \resumeTalk{RFID Skimming 해킹 기법에 관한 고찰}{한국정보보호학회 동계학술대회, 2017}
      \resumeItemListStart
        \resumeItem{김건우, 김태현, 성유원, 여상준, 이진우, 이창엽, 이달원}
      \resumeItemListEnd
  \resumeSubHeadingListEnd


%-----------CERTIFICATIONS-----------------
\section{자격 및 취득 이력}
  \resumeSubHeadingListStart
    \resumeItem{AWS Certified Solutions Architect - Associate (2023.07.26 -- 2026.07.26, 만료)}
    \resumeItem{디지털포렌식 2급 (2022.12.22 -- 2028.12.22)}
    \resumeItem{정보기기운용기능사 (2022.12 취득)}
    \resumeItem{Exterro ACE / AccessData (2022.01.11 -- 2024.01.11, 만료)}
    \resumeItem{인터넷정보관리사 2급 (2020.10 취득)}
    \resumeItem{리눅스마스터 2급 (2019.10 취득)}
    \resumeItem{네트워크관리사 2급 (2016.12 -- 2031.12.13)}
  \resumeSubHeadingListEnd


\end{document}
